\input{papers/template/style.tex}

\begin{document}
  \title{A Theoretical Framework for Quantifying Skill and Luck in Agent Outcomes}

  \author{Giacinto Paolo Saggese$^{*}$}
  \thanks{$^{*}$ University of Maryland}

  \date{\today}

  \maketitle

  \begin{abstract}
    The distinction between skill and luck in determining outcomes represents a fundamental problem in understanding individual and societal achievement. This paper develops a mathematical framework for quantifying the role of luck in agent outcomes by decomposing events into probability, utility, and control components. We formalize luck as value-weighted surprise of outcomes not explained by controllable factors, incorporating concepts from probability theory, decision theory, and causal inference. The framework provides a principled method for rating individual events, measuring aggregate luck distributions across populations, and analyzing societal fairness. Applications span moral philosophy, policy design, and institutional evaluation. The theoretical approach enables empirical measurement of luck's contribution to observed outcomes, supporting evidence-based discussions of merit, responsibility, and distributive justice.
  \end{abstract}

  \setcounter{tocdepth}{2}
  \tableofcontents

  % ############################################################################
  \section{Introduction}

  The attribution of outcomes to skill versus luck influences judgments across domains including moral philosophy, economic policy, educational assessment, and institutional design. Despite widespread intuition that both factors matter, formal methods for separating and quantifying these contributions remain underdeveloped. This gap impedes rigorous analysis of questions such as: To what extent do observed inequalities reflect differences in ability versus circumstance? How should institutions account for luck when evaluating performance? What policies are justified when outcomes depend substantially on factors beyond individual control?

  Existing approaches often conflate luck with randomness or treat it as a residual category. We propose that luck can be formally defined and measured through a multi-component framework that integrates probability (rarity of events), utility (value to the agent), and control (degree of agency over outcomes). This formalization enables quantitative analysis while preserving philosophical distinctions between different forms of luck.

  The framework addresses three interconnected problems. First, it provides a method for rating individual events on a luck scale, accounting for context-dependent information and agent-specific values. Second, it enables decomposition of outcomes into skill and luck components through counterfactual reasoning and expected value calculations. Third, it supports societal-level analysis by aggregating individual luck scores and measuring distributional properties such as luck inequality and the correlation between circumstances and outcomes.

  This paper proceeds as follows. Section~\ref{sec:background} discusses the conceptual foundations and motivation. Section~\ref{sec:theory} develops the mathematical theory, including axioms, probability models, and the core luck function. Section~\ref{sec:rating} presents practical methods for rating events and decomposing skill from luck. Section~\ref{sec:society} extends the framework to societal modeling. Section~\ref{sec:implications} examines consequences for policy and institutions. Section~\ref{sec:conclusion} concludes.

  % ############################################################################
  \section{Background and Motivation}
  \label{sec:background}

  \subsection{Conceptual Foundations}

  The concept of luck appears in diverse contexts with varying definitions. In ordinary language, luck refers to outcomes influenced by chance. In moral philosophy, particularly the literature on moral luck, it denotes factors outside an agent's control that affect judgments of praise or blame. In statistical analysis, luck often represents deviations from expected performance or unexplained variance.

  We propose a unifying definition: luck is the value-weighted surprise of an outcome that the agent did not significantly control, relative to the agent's prior information. This definition incorporates three essential components:

  \begin{itemize}
    \item \textbf{Probability:} Low-probability events contribute more to luck than expected occurrences.
    \item \textbf{Value:} Events must matter to the agent; neutral outcomes are not lucky.
    \item \textbf{Control:} Outcomes predominantly caused by the agent's choices reflect skill rather than luck.
  \end{itemize}

  These components distinguish luck from related concepts. Pure randomness becomes luck only when it affects valued outcomes. Skill represents the portion of outcomes explained by controllable actions. Circumstance refers to background conditions that shape probabilities and opportunities.

  \subsection{Motivation for Formalization}

  Formalizing luck serves multiple purposes. In moral and political philosophy, it clarifies debates about desert and responsibility. If outcomes depend substantially on luck, claims that individuals deserve their positions become more difficult to sustain. In institutional design, understanding luck enables fairer evaluation systems that account for factors beyond individual control. In policy analysis, measuring the role of luck informs debates about taxation, social insurance, and equality of opportunity.

  Empirical measurement requires precise definitions. Without formal structure, assertions about the importance of luck remain speculative. A mathematical framework enables testable predictions, quantitative comparisons across contexts, and evidence-based policy recommendations.

  \subsection{Related Concepts}

  Several mathematical and philosophical traditions inform this framework:

  \begin{itemize}
    \item \textbf{Probability Theory:} Provides the foundation for quantifying surprise and rare events through probability measures and information theory.
    \item \textbf{Decision Theory:} Supplies utility functions for valuing outcomes and expected utility calculations for separating realized outcomes from expectations.
    \item \textbf{Causal Inference:} Offers tools for identifying controllable versus uncontrollable factors through causal graphs and counterfactual reasoning.
    \item \textbf{Game Theory:} Distinguishes strategic skill from chance elements in mixed games of skill and luck.
    \item \textbf{Moral Philosophy:} Examines the normative implications of luck through analyses of moral luck and distributive justice.
  \end{itemize}

  % ############################################################################
  \section{Mathematical Framework}
  \label{sec:theory}

  \subsection{Notation and Definitions}

  Let $\Omega$ denote a sample space of possible events, $\mathcal{F}$ a $\sigma$-algebra on $\Omega$, and $P$ a probability measure. An agent $A$ at time $t$ possesses an information set $I_A(t) \subseteq \mathcal{F}$ representing their knowledge. An event $E \in \mathcal{F}$ occurs with conditional probability $P(E \mid I_A(t))$.

  \begin{definition}[Agent Utility]
    A utility function $U_A: \mathcal{F} \to \mathbb{R}$ assigns value to events from the perspective of agent $A$. Positive values represent favorable outcomes, negative values represent unfavorable outcomes, and zero represents neutral events.
  \end{definition}

  \begin{definition}[Control]
    The control function $C_A: \mathcal{F} \to [0,1]$ measures the degree to which agent $A$ causally influences event $E$. $C_A(E) = 0$ indicates no control (pure chance), while $C_A(E) = 1$ indicates complete control (fully determined by agent's actions).
  \end{definition}

  \subsection{Core Luck Function}

  We define luck as a function of probability, utility, and control.

  \begin{definition}[Luck]
    \label{def:luck}
    The luck experienced by agent $A$ from event $E$ is:
    \begin{equation}
      L_A(E) = U_A(E) \cdot S(E \mid I_A) \cdot (1 - C_A(E))
    \end{equation}
    where $S(E \mid I_A)$ is the surprise function measuring the unexpectedness of $E$ given information $I_A$.
  \end{definition}

  The surprise function $S$ quantifies rarity. A natural choice based on information theory is:
  \begin{equation}
    S(E \mid I_A) = -\log P(E \mid I_A)
  \end{equation}

  This formulation satisfies several desirable properties:

  \begin{proposition}
    The luck function $L_A(E)$ is:
    \begin{enumerate}
      \item Monotonically increasing in outcome value for favorable events
      \item Monotonically decreasing in probability (rarer events are luckier)
      \item Monotonically decreasing in control (less controllable events are luckier)
      \item Zero when $U_A(E) = 0$ (neutral outcomes)
      \item Zero when $C_A(E) = 1$ (fully controlled outcomes)
      \item Information-dependent through $P(E \mid I_A)$
    \end{enumerate}
  \end{proposition}

  \subsection{Alternative Formulations}

  Several variants of the core luck function address different analytical needs.

  \subsubsection{Deviation-Based Luck}

  When repeated trials or comparable scenarios exist, luck can be measured as deviation from expectation:
  \begin{equation}
    L_A^{\text{dev}}(E) = U_A(E) - \mathbb{E}[U_A(E') \mid S_A]
  \end{equation}
  where $S_A$ represents the agent's strategy or controllable choices, and the expectation is taken over events $E'$ that could occur under that strategy. This formulation is particularly useful in skill-luck decomposition for repeated performance evaluation.

  \subsubsection{Fragility-Based Luck}

  Outcomes that would change under small perturbations reflect greater luck:
  \begin{equation}
    L_A^{\text{frag}}(E) = U_A(E) \cdot \left(1 - P(E \mid \text{perturbations})\right)
  \end{equation}
  This captures near-miss scenarios where slight differences in circumstances would have produced substantially different outcomes.

  \subsubsection{Time-Aggregated Luck}

  For sequences of events over time $t = 1, \ldots, T$:
  \begin{equation}
    L_A^{T} = \sum_{t=1}^{T} \delta^t L_A(E_t)
  \end{equation}
  where $\delta \in (0,1]$ is a temporal discount factor. This formulation addresses cumulative luck and path dependence in life trajectories.

  \subsection{Skill-Luck Decomposition}

  Observed outcomes $Y$ can be decomposed into skill and luck components. Let $S_A$ denote the agent's strategy space and $R$ denote random variables outside the agent's control.

  \begin{definition}[Expected Performance]
    The skill-based expected outcome is:
    \begin{equation}
      \bar{Y}_A = \mathbb{E}[Y \mid S_A, I_A]
    \end{equation}
  \end{definition}

  \begin{definition}[Luck Component]
    The luck component of outcome $Y$ is:
    \begin{equation}
      L_A(Y) = Y - \bar{Y}_A
    \end{equation}
  \end{definition}

  This decomposition enables performance evaluation that accounts for uncontrollable factors. In contexts with repeated observations, the variance of outcomes can be partitioned:
  \begin{equation}
    \text{Var}(Y) = \text{Var}(\bar{Y}_A) + \text{Var}(L_A(Y))
  \end{equation}
  where $\text{Var}(\bar{Y}_A)$ represents skill-based variation and $\text{Var}(L_A(Y))$ represents luck-based variation.

  \subsection{Axioms and Properties}

  The framework rests on several axioms that formalize intuitive properties of luck:

  \begin{assumption}[Neutrality]
    Outcomes with zero utility contribute zero luck: $U_A(E) = 0 \implies L_A(E) = 0$.
  \end{assumption}

  \begin{assumption}[Control Dominance]
    Fully controlled outcomes are not lucky: $C_A(E) = 1 \implies L_A(E) = 0$.
  \end{assumption}

  \begin{assumption}[Information Dependence]
    Luck depends on the agent's prior information: $L_A(E)$ is a function of $P(E \mid I_A)$ rather than $P(E)$ alone.
  \end{assumption}

  \begin{assumption}[Value Monotonicity]
    For events with equal probability and control, luck increases with absolute utility: if $P(E_1 \mid I_A) = P(E_2 \mid I_A)$ and $C_A(E_1) = C_A(E_2)$, then $|L_A(E_1)| < |L_A(E_2)|$ when $|U_A(E_1)| < |U_A(E_2)|$.
  \end{assumption}

  % ############################################################################
  \section{Event Rating and Measurement}
  \label{sec:rating}

  \subsection{Rating Pipeline}

  Practical application of the framework requires a systematic procedure for rating events. We propose a step-by-step pipeline:

  \subsubsection{Step 1: Specify Agent and Event}

  Clearly identify the agent $A$, the event $E$, and the time $t$ at which the event occurs. Luck is agent-relative and context-dependent.

  \subsubsection{Step 2: Define Reference Model}

  Construct a probability model $P(\cdot \mid I_A(t))$ representing the agent's information state prior to the event. This may involve:
  \begin{itemize}
    \item Historical base rates for similar events
    \item Statistical models fitted to relevant data
    \item Expert forecasts or prediction markets
    \item Agent-specific models accounting for their knowledge and position
  \end{itemize}

  \subsubsection{Step 3: Compute Surprise}

  Calculate the surprise score:
  \begin{equation}
    S(E) = -\log P(E \mid I_A)
  \end{equation}
  Low-probability events receive higher surprise scores.

  \subsubsection{Step 4: Assess Outcome Value}

  Determine utility $U_A(E)$ using appropriate scales:
  \begin{itemize}
    \item Monetary value for financial outcomes
    \item Health metrics (life expectancy, quality-adjusted life years) for health outcomes
    \item Career advancement measures for professional outcomes
    \item Normalized scores for comparative analysis
  \end{itemize}

  \subsubsection{Step 5: Estimate Control}

  Quantify the degree of agent control $C_A(E) \in [0,1]$. Methods include:
  \begin{itemize}
    \item Causal modeling to identify controllable versus uncontrollable factors
    \item Variance decomposition showing proportion explained by agent's actions
    \item Counterfactual analysis examining outcome sensitivity to agent's choices
    \item Repeatability analysis measuring consistency of outcomes across similar scenarios
  \end{itemize}

  \subsubsection{Step 6: Calculate Luck Score}

  Apply the luck function:
  \begin{equation}
    L_A(E) = U_A(E) \cdot S(E) \cdot (1 - C_A(E))
  \end{equation}

  \subsubsection{Step 7: Normalize}

  For cross-event or cross-agent comparison, normalize within relevant populations:
  \begin{equation}
    L_A^*(E) = \frac{L_A(E) - \mu_L}{\sigma_L}
  \end{equation}
  producing a standardized luck z-score.

  \subsection{Data Requirements}

  Empirical application requires several types of data:

  \begin{table}[H]
    \centering
    \begin{tabular}{ll}
      \toprule
      \textbf{Data Type} & \textbf{Purpose} \\
      \midrule
      Probability data & Estimate $P(E \mid I_A)$ \\
      Historical frequencies & Base rates for similar events \\
      Information records & Define information set $I_A$ \\
      Outcome measurements & Quantify utility $U_A(E)$ \\
      Decision logs & Track controllable actions \\
      Causal data & Estimate control $C_A(E)$ \\
      Baseline performance & Compute expected values \\
      Population distributions & Normalize luck scores \\
      \bottomrule
    \end{tabular}
    \caption{Data requirements for empirical luck measurement}
    \label{tab:data}
  \end{table}

  \subsection{Example Application}

  Consider a concrete example: an entrepreneur's startup succeeds after raising venture capital. We analyze the luck component:

  \begin{itemize}
    \item \textbf{Event:} Startup achieves successful exit within five years
    \item \textbf{Probability:} Base rate for similar ventures is approximately 10\%, so $P(E) = 0.1$ and $S(E) = -\log(0.1) \approx 2.3$
    \item \textbf{Utility:} Financial gain of \$5M, normalized to $U_A(E) = 5$
    \item \textbf{Control:} Entrepreneur's strategy and execution contributed substantially, but market conditions and network effects were also crucial; estimate $C_A(E) = 0.6$
    \item \textbf{Luck Score:} $L_A(E) = 5 \times 2.3 \times (1 - 0.6) = 4.6$
  \end{itemize}

  This positive luck score indicates that while the entrepreneur's skill was important, favorable circumstances and timing contributed significantly to the outcome.

  % ############################################################################
  \section{Societal Modeling and Aggregate Measures}
  \label{sec:society}

  \subsection{Modeling Life Outcomes}

  To analyze luck at the societal level, we model individual life outcomes as a function of three input categories:

  \begin{itemize}
    \item \textbf{Circumstances} $B_i$: Factors not chosen by the individual (parental background, birth location, cohort, genetics)
    \item \textbf{Choices} $A_i$: Factors under individual control (effort, decisions, strategies)
    \item \textbf{Shocks} $S_i$: Random events affecting the individual (accidents, encounters, policy changes)
  \end{itemize}

  A general outcome model takes the form:
  \begin{equation}
    Y_i = f(B_i, A_i) + g(S_i) + \varepsilon_i
  \end{equation}
  where $f$ represents the deterministic relationship between circumstances, choices, and outcomes, $g$ captures the effect of random shocks, and $\varepsilon_i$ represents measurement error or small unmodeled factors.

  \subsection{Individual Luck Measures}

  Three complementary measures capture different aspects of luck:

  \subsubsection{Residual Luck}

  Residual luck measures the deviation of actual outcomes from predictions based on circumstances and choices:
  \begin{equation}
    L_i^{\text{res}} = Y_i - \mathbb{E}[Y_i \mid B_i, A_i]
  \end{equation}
  This represents the core luck component—the portion of outcomes not explained by observable inputs.

  \subsubsection{Surprise-Weighted Luck}

  For outcomes far from expectations, surprise weighting emphasizes extreme cases:
  \begin{equation}
    L_i^{\text{sur}} = U(Y_i) \cdot \left[-\log P(Y_i \mid B_i, A_i)\right]
  \end{equation}

  \subsubsection{Control-Adjusted Luck}

  When control varies across individuals or contexts:
  \begin{equation}
    L_i^{\text{ctrl}} = (Y_i - \mathbb{E}[Y_i \mid B_i, A_i]) \cdot (1 - C_i)
  \end{equation}
  where $C_i$ represents the degree of control individual $i$ has over their outcomes.

  \subsection{Aggregate Luck Metrics}

  Societal-level analysis aggregates individual luck scores to characterize distributional properties.

  \subsubsection{Luck Inequality}

  The Gini coefficient of absolute luck values measures dispersion:
  \begin{equation}
    G_L = \frac{\sum_{i=1}^{n} \sum_{j=1}^{n} |L_i - L_j|}{2n^2\bar{L}}
  \end{equation}
  High luck inequality indicates that some individuals experience far more favorable or unfavorable luck than others.

  \subsubsection{Circumstantial Dependence}

  The proportion of outcome variance explained by circumstances (rather than choices or luck) measures the degree to which outcomes are predetermined by birth:
  \begin{equation}
    \rho_B = \frac{\text{Var}(\mathbb{E}[Y \mid B])}{\text{Var}(Y)}
  \end{equation}
  Higher values indicate less social mobility and greater importance of circumstantial luck.

  \subsubsection{Agency Contribution}

  The additional variance explained by choices beyond circumstances:
  \begin{equation}
    \rho_A = \frac{\text{Var}(\mathbb{E}[Y \mid B, A]) - \text{Var}(\mathbb{E}[Y \mid B])}{\text{Var}(Y)}
  \end{equation}
  This measures the role of individual agency in determining outcomes.

  \subsubsection{Mobility Index}

  Intergenerational mobility can be measured through the correlation between parent and child outcomes:
  \begin{equation}
    M = 1 - \text{Corr}(Y_{\text{parent}}, Y_{\text{child}})
  \end{equation}
  Lower correlation (higher $M$) suggests that circumstances of birth have less influence on eventual outcomes.

  \subsection{Empirical Estimation}

  Practical estimation typically proceeds through regression-based variance decomposition. Consider a model:
  \begin{equation}
    Y_i = \alpha + \beta_B B_i + \beta_A A_i + \varepsilon_i
  \end{equation}

  Stepwise regression allows decomposition:
  \begin{enumerate}
    \item Regress $Y$ on $B$ alone: variance explained is $R^2_B$
    \item Regress $Y$ on both $B$ and $A$: variance explained is $R^2_{B,A}$
    \item Residual variance: $1 - R^2_{B,A}$ represents luck
    \item Agency contribution: $R^2_{B,A} - R^2_B$
  \end{enumerate}

  % ############################################################################
  \section{Implications and Applications}
  \label{sec:implications}

  \subsection{Moral and Philosophical Implications}

  The framework clarifies debates about desert and responsibility. If formal analysis demonstrates that outcomes depend substantially on luck rather than controllable factors, several philosophical implications follow:

  \begin{itemize}
    \item \textbf{Desert:} Claims that individuals fully deserve their outcomes become difficult to sustain when luck plays a significant role.
    \item \textbf{Responsibility:} Moral evaluation must account for the limited control agents have over their circumstances and the events that befall them.
    \item \textbf{Dignity:} If outcomes are substantially luck-determined, human worth should not be closely tied to achievement.
  \end{itemize}

  These implications align with philosophical analyses of moral luck, which argue that factors beyond an agent's control affect our moral judgments inappropriately.

  \subsection{Policy Implications}

  Quantifying the role of luck informs policy debates across multiple domains.

  \subsubsection{Taxation and Redistribution}

  When luck contributes substantially to income and wealth, the case for progressive taxation and redistribution strengthens. If high earners benefit from favorable circumstances and chance events, claims of full entitlement to their income weaken. Luck-adjusted taxation could incorporate measurements of the luck component in individual outcomes.

  \subsubsection{Social Insurance}

  Recognition that adverse outcomes often reflect bad luck rather than poor choices supports comprehensive social insurance. Unemployment benefits, healthcare access, and disability support become justified as protections against uncontrollable risks rather than rewards for underperformance.

  \subsubsection{Educational Policy}

  Measuring luck in educational outcomes can inform admission policies and resource allocation. If standardized test performance reflects substantial circumstantial luck (parental resources, school quality, neighborhood effects), policies that emphasize fine-grained score differences become less defensible. Alternative approaches such as lottery-based admission among qualified candidates warrant consideration.

  \subsection{Institutional Design}

  Organizations can incorporate luck awareness in evaluation and reward systems.

  \subsubsection{Performance Evaluation}

  Corporations and public institutions that evaluate employees on outcomes without accounting for luck risk misattributing results. Luck-adjusted performance metrics provide fairer assessments by comparing actual outcomes to expected outcomes given circumstances and constraints.

  \subsubsection{Promotion and Selection}

  When selecting among candidates with similar qualifications, acknowledging that fine-grained differences may reflect luck rather than ability supports lottery-based selection or increased weight on diverse criteria beyond narrow performance measures.

  \subsubsection{Executive Compensation}

  CEO pay often reflects firm performance that depends substantially on market conditions and timing. Luck-adjusted compensation schemes could separate controllable performance from market-driven outcomes, reducing excessive rewards for favorable circumstances.

  \subsection{Cultural and Psychological Effects}

  Wider recognition of luck's role may produce several cultural shifts:

  \begin{itemize}
    \item \textbf{Increased Humility:} Successful individuals may develop greater appreciation for favorable circumstances rather than attributing outcomes entirely to personal merit.
    \item \textbf{Reduced Stigma:} Failure may be viewed with less moral judgment when understood as partially reflecting bad luck.
    \item \textbf{Systems Thinking:} Cultural narratives may shift from individual hero stories toward recognition of structural conditions and systemic factors.
    \item \textbf{Risk of Fatalism:} Overemphasis on luck without careful communication could reduce motivation and effort, particularly in educational contexts.
  \end{itemize}

  \subsection{Limitations and Caveats}

  Several limitations warrant acknowledgment. First, quantifying control remains challenging and often requires judgment calls. Second, the framework assumes agents have definable utility functions, which may not capture all relevant dimensions of value. Third, practical application requires substantial data that may not always be available. Fourth, emphasizing luck could produce unintended psychological effects if not communicated carefully.

  Despite these limitations, the framework provides a structured approach to a question that otherwise remains informal and speculative.

  % ############################################################################
  \section{Conclusion}
  \label{sec:conclusion}

  This paper develops a formal framework for quantifying skill and luck in agent outcomes. By integrating concepts from probability theory, decision theory, and causal inference, the approach provides mathematically precise definitions and measurement procedures. The framework distinguishes three essential components—probability, utility, and control—and combines them in a multiplicative luck function that satisfies intuitive axioms.

  Applications span individual event rating, performance evaluation, and societal-level analysis. At the individual level, the framework enables luck-adjusted assessment of outcomes. At the aggregate level, it supports measurement of luck inequality, circumstantial dependence, and social mobility. These measurements inform debates about fairness, desert, and policy design.

  Several directions for future work appear promising. Empirical applications to specific domains such as labor markets, financial markets, or educational outcomes could validate the framework and refine estimation methods. Extensions incorporating dynamic considerations such as path dependence and cumulative advantage would enhance understanding of lifetime trajectories. Integration with behavioral economics and psychology could address perception biases and motivation effects. Normative analysis of optimal policy under measured luck distributions could translate descriptive findings into prescriptive recommendations.

  The framework suggests that outcomes depend more substantially on luck than commonly acknowledged. This recognition has implications for how societies structure incentives, allocate resources, and make judgments about individual responsibility. By providing formal tools for measurement and analysis, the framework enables more rigorous and evidence-based engagement with these questions.

  %\bibliographystyle{amsrn}

  \begin{thebibliography}{99}

  \bibitem{williams1981}
  Williams, B. (1981). \textit{Moral Luck}. Cambridge University Press.

  \bibitem{nagel1979}
  Nagel, T. (1979). Moral luck. In \textit{Mortal Questions} (pp. 24--38). Cambridge University Press.

  \bibitem{rawls1971}
  Rawls, J. (1971). \textit{A Theory of Justice}. Harvard University Press.

  \bibitem{kahneman2011}
  Kahneman, D. (2011). \textit{Thinking, Fast and Slow}. Farrar, Straus and Giroux.

  \bibitem{mauboussin2012}
  Mauboussin, M. J. (2012). \textit{The Success Equation: Untangling Skill and Luck in Business, Sports, and Investing}. Harvard Business Review Press.

  \bibitem{frank2016}
  Frank, R. H. (2016). \textit{Success and Luck: Good Fortune and the Myth of Meritocracy}. Princeton University Press.

  \bibitem{pearl2009}
  Pearl, J. (2009). \textit{Causality: Models, Reasoning, and Inference} (2nd ed.). Cambridge University Press.

  \bibitem{shannon1948}
  Shannon, C. E. (1948). A mathematical theory of communication. \textit{Bell System Technical Journal}, 27(3), 379--423.

  \bibitem{savage1954}
  Savage, L. J. (1954). \textit{The Foundations of Statistics}. Wiley.

  \bibitem{roemer1998}
  Roemer, J. E. (1998). \textit{Equality of Opportunity}. Harvard University Press.

  \end{thebibliography}

  % Original bibliography command (commented out):
  % \bibliography{references}
\end{document}
